On a reproduit l'étiquette d'une solution commerciale d'acide chlorhydrique.

\begin{minipage}{\linewidth}
    \begin{figure}[H]
        \centering
        \includegraphics[width=\textwidth]{2025-12-09_000006-}
    \end{figure}
\end{minipage}

$35\%$ est le titre massique de cette solution. Cela signifie que dans
$\qty{100}{\g}$ de solution, on a dissous $\qty{35}{\g}$ de chlorure
d'hydrogène.
\begin{enumerate}
    \item Quelles précautions doit-on prendre pour manipuler cette solution~?
    \item Calculer sa concentration molaire.
    \item La capacité du flacon est de $\qty{5.0}{\L}$. Quel volume de chlorure
          d'hydrogène a-t-on dissous, dans les conditions où le volume molaire
          vaut $\qty{24.0}{\L\per\mol}$ dans ce flacon~?
    \item Cette solution étant trop concentrée, on désire la diluer au
          $\frac{1}{10}$~; proposer un mode opératoire respectant les consignes
          de sécurité pour obtenir $\qty{100}{\mL}$ de solution diluée.
\end{enumerate}


